\section{Hardware e sensores}
\label{sec:hardware}

O nó sensor usa uma placa de desenvolvimento com o ESP32, um SoC com dois núcleos Xtensa, Wi-Fi integrado, conversores analógico-digitais (ADC) do tipo SAR, controladores UART e periférico de toque capacitivo \citep{esp32datasheet,esp32trm}. A Tabela~\ref{tab:pinagem} resume a ligação dos sensores. \pendente{Informar o modelo exato da placa de desenvolvimento (ex.: DevKit V1 / DevKitC) e incluir foto ou esquemático da montagem.}

\begin{table*}[!t]
\caption{Pinagem dos sensores no protótipo. GPIO34 e GPIO35 são pinos somente de entrada no ESP32 e não possuem resistores internos de \textit{pull-up}/\textit{pull-down} \citep{esp32datasheet}.}
\label{tab:pinagem}
\centering
\small
\begin{tblr}{
  colspec={X[1.2,l] X[1.1,l] X[0.8,l] X[1.6,l] X[2.2,l]},
  row{1}={font=\bfseries},
  hline{1,Z}={0.8pt}, hline{2}={0.4pt},
  rowsep=1pt
}
Sensor & Pino do módulo & ESP32 & Periférico / configuração & Observação \\
Microfone (LM393) & AO & GPIO34 & ADC1, canal 6; modo contínuo a 20\,kHz, atenuação de 12\,dB & módulo alimentado em 3,3\,V para que AO não ultrapasse a faixa do ADC \\
Microfone (LM393) & DO & GPIO35 & entrada digital com interrupção em qualquer borda & limiar ajustado no trimpot do módulo; o \textit{pull-up} é o do próprio módulo \\
Microfone (LM393) & VCC / GND & 3V3 / GND & alimentação & --- \\
Sonar LV-MaxSonar-EZ1 & TX (pino 5) & GPIO16 & UART2 RX, 9600\,bit/s, 8N1, linha RX invertida & saída serial com polaridade RS-232 e níveis 0--Vcc \\
Sonar LV-MaxSonar-EZ1 & +5 (pino 6) / GND (pino 7) & 3V3 / GND & alimentação (2,5--5,5\,V segundo o fabricante) & em 3,3\,V a saída TX já fica em nível compatível com o ESP32 \\
Sonar LV-MaxSonar-EZ1 & RX, AN, PW, BW & --- & desconectados & RX aberto: medição contínua; BW aberto: saída serial \\
Capacitivo CR18-8DP & saída PNP & GPIO25 & entrada digital & exige adaptação de nível se alimentado acima de 3,3\,V \pendente{circuito de interface real} \\
Capacitivo CR18-8DP & V+ / 0\,V & \pendente{fonte} & alimentação externa \pendente{tensão usada} & terra comum com o ESP32 \\
Botão BOOT & --- & GPIO0 & entrada digital & mantido por 3\,s apaga a configuração (NVS) e reinicia \\
\end{tblr}
\end{table*}

\subsection{Microfone de eletreto com comparador LM393}
\label{sec:mic}

\paragraph{Princípio e módulo.}
O módulo usado (placa vermelha com saídas AO e DO) contém uma cápsula de microfone de eletreto polarizada por resistor e um comparador LM393 \citep{tilm393}. A saída analógica AO entrega a tensão da cápsula sobreposta a um nível DC; o comparador compara esse sinal com um limiar ajustado por trimpot e produz a saída digital DO. O módulo \textbf{não possui estágio de amplificação} no caminho de AO, de modo que a amplitude do sinal de voz do gato é pequena em relação à faixa do ADC. \pendente{Medir com osciloscópio a amplitude pico a pico de AO para um miado a 0,5\,m e 1,5\,m e o nível DC de repouso.}

\paragraph{Ligação.}
AO vai ao GPIO34 (ADC1, canal 6). Usa-se o ADC1 porque o ADC2 do ESP32 é compartilhado com o driver de Wi-Fi \citep{espidfadc}. DO vai ao GPIO35, configurado com interrupção em qualquer borda: cada borda incrementa um contador, lido e zerado a cada quadro de análise, o que ajuda a marcar o início do som. O módulo é alimentado em 3,3\,V para que AO permaneça dentro da faixa de entrada do ADC.

\paragraph{Condicionamento digital.}
Como não há filtragem analógica, todo o condicionamento é feito no firmware (\texttt{mic\_sensor.c}):
\begin{enumerate}
  \item \textbf{Amostragem}: o driver contínuo do ADC (\textit{DMA}) é configurado em 20\,kHz, o mínimo aceito nesse modo no ESP32, com atenuação de 12\,dB \citep{espidfadc}.
  \item \textbf{Decimação 2:1}: cada par de amostras é substituído pela média, resultando em $f_s = 10$\,kHz. A média de pares é um filtro FIR simples (passa-baixas com resposta $|\cos(\pi f / 20\,\text{kHz})|$) e também corrige a troca de pares de amostras do ADC do ESP32 em modo contínuo.
  \item \textbf{Remoção do nível DC}: uma média móvel exponencial $d[n] = d[n-1] + \alpha\,(x[n]-d[n-1])$, com $\alpha = 0{,}001$, estima o nível DC, que é subtraído do sinal. Isso equivale a um passa-altas de primeira ordem com frequência de corte aproximada de $\alpha f_s / 2\pi \approx 1{,}6$\,Hz.
  \item \textbf{Quadros}: o sinal é processado em quadros de 256 amostras ($\approx 25{,}6$\,ms).
\end{enumerate}

Para cada quadro calculam-se:
\begin{itemize}
  \item o \textbf{volume relativo}, em dB, sobre as contagens do ADC:
  \begin{equation}
    L = 20\log_{10}\!\left(\sqrt{\tfrac{1}{N}\textstyle\sum_{i} x_i^2} + 10^{-3}\right);
    \label{eq:rms}
  \end{equation}
  \item a \textbf{taxa de cruzamentos por zero} (ZCR), fração de amostras em que o sinal troca de sinal --- alta em ruídos de banda larga como o sibilo;
  \item a \textbf{frequência fundamental} $f_0$, pela autocorrelação normalizada
  \begin{equation}
    r(\tau) = \frac{\sum_i x_i\,x_{i+\tau}}{\sqrt{\sum_i x_i^2 \,\sum_i x_{i+\tau}^2}},
    \label{eq:autocorr}
  \end{equation}
  avaliada para atrasos correspondentes a 80--1500\,Hz. O quadro é considerado \textbf{tonal} se o máximo de $r(\tau)$ for pelo menos 0,5; o período é o primeiro pico com pelo menos 90\,\% do máximo (evitando erro de oitava) e é refinado por interpolação parabólica, $f_0 = f_s/\tau^*$.
\end{itemize}

\paragraph{Detecção de eventos.}
No primeiro segundo (40 quadros) o firmware mede o ruído de fundo. Depois, um quadro é considerado ``alto'' se $L$ supera o ruído de fundo em 12\,dB ou se houve borda em DO; nos quadros silenciosos o ruído de fundo continua se adaptando lentamente. Um som termina após 200\,ms de silêncio; pausas maiores que 40\,ms dentro do som contam como novas rajadas; sons com menos de 150\,ms (estalos) são descartados e sons com mais de 10\,s (TV, aspirador) são cortados. Ao final, o firmware envia a duração, o pico de volume, $f_0$ médio, de início e de fim, a fração de quadros tonais, a ZCR média, o número de rajadas e um envelope de volume com até 64 pontos (máximo de cada intervalo, para preservar os picos).

\paragraph{Limitações.}
\begin{itemize}
  \item \textbf{Volume apenas relativo}: a Equação~\eqref{eq:rms} opera sobre contagens do ADC, sem calibração em nível de pressão sonora; o valor só é comparável entre medições do mesmo módulo, na mesma posição e com o mesmo ganho.
  \item \textbf{Resolução efetiva baixa}: sem pré-amplificador, o sinal ocupa uma fração pequena da faixa do ADC, e o ruído e a não linearidade do ADC do ESP32 \citep{esp32datasheet} pesam proporcionalmente mais.
  \item \textbf{Ausência de filtro \textit{anti-aliasing}}: componentes acima de 10\,kHz se rebatem na amostragem a 20\,kHz, e a média de pares atenua apenas parcialmente a faixa de 5 a 10\,kHz antes da decimação. Para a estimativa de $f_0$ (até 1,5\,kHz) isso é tolerável, mas afeta a ZCR.
  \item \textbf{Alcance e sensibilidade}: a equipe observou alcance útil de cerca de 0,5--1,5\,m e que o ronronar não é captado. \pendente{Confirmar com medições sistemáticas.}
  \item Com o módulo de saída apenas digital (FC-04 ou placa preta), o firmware pode ser compilado com \texttt{MIC\_HAS\_ANALOG = 0}, enviando apenas duração e rajadas.
\end{itemize}

\subsection{Sonar ultrassônico LV-MaxSonar-EZ1}
\label{sec:sonar}

\paragraph{Princípio.}
O LV-MaxSonar-EZ1 (MB1010) emite pulsos de 42\,kHz e mede o tempo de voo do eco. Segundo o fabricante \citep{maxbotixez1}, opera de 2,5 a 5,5\,V, fornece leituras a cada 49\,ms ($\approx 20$\,Hz) com resolução de 1 polegada entre 6 e 254 polegadas, e \textbf{objetos entre 0 e 6 polegadas (15,24\,cm) são reportados como 6 polegadas}. As três saídas (largura de pulso, 147\,$\mu$s/pol; tensão analógica, Vcc/512 por polegada; e serial) ficam ativas simultaneamente. Na partida, o sensor executa um ciclo de calibração do padrão de \textit{ringdown}; o fabricante recomenda que não haja objetos a menos de 7 polegadas (idealmente 14) durante esse ciclo, e informa que, operando em 3,3\,V, as distâncias tendem a ser reportadas 1 a 2\,\% maiores.

\paragraph{Ligação e leitura.}
Usa-se a saída serial: o pino TX envia, a 9600\,bit/s (8N1), um caractere \texttt{R}, três dígitos ASCII com a distância em polegadas e um \textit{carriage return}. Os níveis são 0--Vcc, mas com a polaridade do RS-232 (lógica invertida) \citep{maxbotixez1}. Por isso o firmware configura a UART2 do ESP32 \citep{espidfuart} com a linha RX invertida em \textit{hardware} (\texttt{uart\_set\_line\_inverse}), dispensando um inversor externo. Com o sonar alimentado em 3,3\,V, a saída já é compatível com a entrada do ESP32. O analisador de pacotes aceita somente o formato \texttt{R} + 3 dígitos + \texttt{CR} e converte polegadas em centímetros (resolução de 2,54\,cm).

\paragraph{Detecção de uso da caixa de areia.}
O sonar é instalado acima da caixa, apontado para baixo. A detecção usa uma linha de base (\textit{baseline}) e histerese:
\begin{enumerate}
  \item nas 20 primeiras leituras ($\approx 1$\,s) a linha de base é a média das distâncias (caixa vazia);
  \item a \textbf{entrada} é confirmada quando a distância fica pelo menos 10\,cm abaixo da linha de base por 10 leituras seguidas ($\approx 0{,}5$\,s);
  \item a \textbf{saída} é confirmada quando a queda volta a no máximo 6\,cm por 30 leituras seguidas ($\approx 1{,}5$\,s);
  \item com a caixa livre e variação menor que 6\,cm, a linha de base acompanha derivas lentas por média móvel exponencial (peso 0,02 por amostra).
\end{enumerate}
Os limiares diferentes de entrada e saída funcionam como a histerese de um comparador Schmitt, evitando oscilação do estado quando a distância fica próxima do limiar. A entrada gera a leitura 1 e a saída, a leitura 0.

\paragraph{Observação no protótipo e limitação importante.}
No registro serial da placa, a linha de base calibrada foi de 17,8\,cm, valor que corresponde a exatamente 7 polegadas. Como o sensor não reporta valores abaixo de 6 polegadas (15,24\,cm), a maior queda mensurável nessa montagem é de cerca de 2,5\,cm, inferior ao limiar de entrada de 10\,cm: \textbf{nessa geometria a entrada do gato não pode ser detectada}. Além disso, 7 polegadas é o mínimo de distância livre recomendado na calibração de partida. Recomenda-se instalar o sonar de forma que a superfície da areia fique a pelo menos $\approx 26$\,cm (15,24\,cm + 10\,cm + margem) da face do sensor, ou rever o limiar. \pendente{Verificar a altura real de instalação, se a leitura de 17,8\,cm era da areia ou de outro objeto (borda, tampa), e repetir o teste após o ajuste.}

\subsection{Sensor capacitivo de proximidade CR18-8DP}
\label{sec:capacitivo}

\paragraph{Princípio.}
Sensores capacitivos de proximidade industriais possuem um oscilador cujo eletrodo, na face do sensor, forma uma capacitância com o meio à frente. Quando um material com permissividade maior que a do ar (como a água) se aproxima, a capacitância aumenta, a amplitude de oscilação muda e um comparador interno aciona o transistor de saída. Assim, o sensor detecta o conteúdo do pote através de uma parede plástica fina. O CR18-8DP é um sensor cilíndrico M18 com saída digital PNP (liga/desliga); pela nomenclatura usual do fabricante, ``8'' indica a distância nominal de detecção em milímetros e ``DP'', alimentação DC com saída PNP. \pendente{Confirmar no rótulo/folha de dados: faixa de tensão de alimentação, tipo de contato (NA/NF), distância de detecção e existência de trimpot de sensibilidade.}

\paragraph{Ligação e leitura.}
A saída é lida no GPIO25 por um módulo \texttt{capacitive\_sensor} do firmware ESP-IDF, que envia 0\,\% ou 100\,\%. Esse módulo está presente no protótipo (compilação local), mas \textbf{ainda não foi enviado ao repositório}. \pendente{Publicar o código do módulo e descrever sua lógica (amostragem, filtragem de \textit{bounce}, periodicidade do envio).}

\paragraph{Adaptação de nível.}
Sensores desse tipo costumam ser alimentados entre 10 e 30\,V. Com saída PNP, quando acionada a saída fica próxima de V+, tensão incompatível com o GPIO do ESP32, que opera em 3,3\,V \citep{esp32datasheet}; quando não acionada, a saída fica em aberto. É necessário, portanto, um circuito de interface, por exemplo:
\begin{itemize}
  \item \textbf{divisor resistivo} da saída para o terra, que também serve de \textit{pull-down}. Para V+ = 12\,V, $R_1 = 10$\,k$\Omega$ e $R_2 = 3{,}3$\,k$\Omega$ resultam em $12 \cdot 3{,}3/13{,}3 \approx 2{,}98$\,V no GPIO;
  \item \textbf{optoacoplador}, que isola galvanicamente a fonte do sensor do ESP32;
  \item \textbf{transistor NPN} em emissor comum com \textit{pull-up} para 3,3\,V (lógica invertida).
\end{itemize}
\pendente{Descrever o circuito de interface efetivamente montado, com valores e medição da tensão no GPIO25 nos dois estados.}

\paragraph{Limitações.}
A saída indica apenas \textbf{presença ou ausência de material na altura em que o sensor está instalado}, e não um nível contínuo: o valor 100\,\% significa ``há água/ração na altura do sensor''. Para níveis discretos seriam necessários vários sensores em alturas diferentes. A ração seca tem permissividade e densidade bem menores que a água, o que reduz a variação de capacitância e pode exigir ajuste de sensibilidade; a aproximação do próprio gato também pode acionar o sensor. \pendente{Registrar testes com pote cheio/vazio, com água e com ração, e com o gato próximo.}

\paragraph{Abordagem anterior.}
O esboço inicial em Arduino (\texttt{esp32-petcare.ino}) previa medir o nível pelo periférico de toque capacitivo do próprio ESP32 (canais T0/GPIO4 e T3/GPIO15), que conta ciclos de carga e descarga do eletrodo e cuja leitura diminui quando a capacitância aumenta \citep{espidftouch}, com mapeamento linear entre leituras de pote vazio e cheio. Essa abordagem não foi levada ao firmware ESP-IDF. \pendente{Explicar por que se optou pelo CR18-8DP (sensibilidade, robustez, disponibilidade).}
